\chapter{초록}

본 보고서는 서울대학교 인문대학 언어학과의 전공 선택 과목인 `언어조사 및 분석'을 수강한 학생들에 의해 작성되었다. 해당 보고서의 내용은 1박 2일에 걸쳐 조사한 벌교 방언의 채록과 관련하여 사전에 준비한 것, 사후에 자료를 정리한 것, 그리고 채록된 자료 그 자체를 알아보기 쉽도록 정리한 것이 주를 이룬다.

조사 대상 방언은 벌교 방언이었으며, 해당 방언이 선정된 것은 현재 빠른 속도로 소멸하고 있는 서남 방언의 일부로서 그 고유한 특징을 보존하고 채록할 가치가 있다고 믿어지기 때문이었다. 이에 대한 부연은 본문의 `조사 지역 및 언어' 장에 정리되어 있다. 현지 조사는 대한민국 전라남도 보성군 벌교읍에서 진행되었으나, 현지조사 2일차인 11월 1일에는 조사를 진행할 장소가 마땅치 않아, 벌교 방언의 화자를 인근 지역인 전라남도 순천시 별량면으로 모셔온 후에 조사를 진행하였다.

전체 과정은 2025년 9월 2일부터 2025년 12월 22일까지 세 단계로 구분된다. 첫 단계는 언어조사 이론 학습 및 현지 조사 준비에 집중한 조사 준비 단계(9월 2일\textasciitilde{}10월 30일)로,  이 기간에 대한 정보는 `조사 준비' 장에 정리되어 있다. 두 번째 단계와 세 번째 단계는 각각 1박 2일간 현지에서 방언 자료를 채록한 현지 조사 단계(10월 31일\textasciitilde{}11월 1일)와 채록된 자료의 전사, 전산화, 분석 및 최종 보고서 작성이 이루어진 후속 분석 단계(11월 2일\textasciitilde{}12월 22일)로, 이 두 기간에 대한 정보는 `조사 과정 및 결과' 장에 정리되어 있다.

벌교 방언을 채록한 녹음 자료 및 영상 자료에 대한 상세 정보는 부록 B `파일 목록'에 정리되어 있으며, 벌교 방언 채록 자료를 전사한 것은 부록 C `전사문'에 정리되어 있다. 그리고 그것들을 가공하여 벌교 방언의 어휘 체계 및 문법 체계를 드러낸 것은 부록 A `어휘·문법 목록'에 정리되어 있다. 해당 과목에 참여한 인원은 총 9인으로, 교수자 1인, 학생 6인, 그리고 언어조사의 전 과정을 보조한 조교 1인으로 구성된 수업 참여 인원과, 조사 과정 중에 학생 인솔에 참여한 외부 인원 1인으로 이루어졌으며, 학생 6인의 세부적인 인적사항은  `후기 및 사진' 장에 정리되어 있다. 이외에 방언 연구에 참여하는 후속 연구 인원들을 위한 정보 및 자료는 `개선 사항 및 제언' 장 및 `후기 및 사진' 장에 정리되어 있다.
